% ==============================================================================
% DreamDojo: A Foundation World Model for Open-World Dexterous Robot Tasks
% Paper Review Presentation — WoRV / MaumAI
% ==============================================================================

% ==============================================================================
% Paper2PR — Shared Beamer Preamble
% ==============================================================================
% Paper-agnostic preamble for all paper review presentations.
% Colors and environments match Quarto clean-academic.scss theme.
% ==============================================================================

\documentclass[aspectratio=169, 10pt]{beamer}

% --- Theme ---
\usetheme{default}
\usecolortheme{default}
\setbeamertemplate{navigation symbols}{}
\setbeamertemplate{footline}[frame number]
\setbeamertemplate{itemize items}[circle]
\setbeamertemplate{enumerate items}[default]

% --- Colors (matching Quarto clean-academic.scss) ---
\usepackage{xcolor}
\definecolor{primaryblue}{HTML}{012169}
\definecolor{primarygold}{HTML}{B9975B}
\definecolor{highlightyellow}{HTML}{F2A900}
\definecolor{lightbg}{HTML}{E8EDF5}
\definecolor{jet}{HTML}{1A1A1A}
\definecolor{positivegreen}{HTML}{15803D}
\definecolor{negativered}{HTML}{B91C1C}

\setbeamercolor{title}{fg=primaryblue}
\setbeamercolor{subtitle}{fg=primarygold}
\setbeamercolor{frametitle}{fg=primaryblue}
\setbeamercolor{structure}{fg=primaryblue}
\setbeamercolor{normal text}{fg=jet}
\setbeamercolor{itemize item}{fg=primaryblue}
\setbeamercolor{itemize subitem}{fg=primarygold}
\setbeamercolor{enumerate item}{fg=highlightyellow}

% --- Fonts (XeLaTeX) ---
\usepackage{fontspec}
% Source Sans Pro (static OTF from TexLive — explicit path for reliability)
\setsansfont{SourceSansPro}[
  Path = /usr/local/texlive/2024/texmf-dist/fonts/opentype/adobe/sourcesanspro/,
  UprightFont = SourceSansPro-Regular.otf,
  BoldFont = SourceSansPro-Bold.otf,
  ItalicFont = SourceSansPro-RegularIt.otf,
  BoldItalicFont = SourceSansPro-BoldIt.otf,
  Scale = 1.0
]
\usefonttheme{professionalfonts}

% --- Math ---
\usepackage{amsmath}
\usepackage{amssymb}
\usepackage{bm}
\usepackage{mathtools}

% --- Graphics ---
\usepackage{graphicx}
\usepackage{tikz}
\usepackage{booktabs}
\usepackage{multicol}
\usepackage{hyperref}

% --- tcolorbox for custom environments ---
\usepackage[most]{tcolorbox}

% methodbox — Blue-bordered for technical details
\newtcolorbox{methodbox}[1][]{%
  colback=primaryblue!6,
  colframe=primaryblue,
  coltitle=white,
  fonttitle=\bfseries,
  left=8pt, right=8pt, top=6pt, bottom=6pt,
  boxrule=0pt, leftrule=4pt,
  arc=4pt,
  title={#1}
}

% keybox — Gold-bordered for key insights
\newtcolorbox{keybox}[1][]{%
  colback=primarygold!12,
  colframe=primarygold,
  coltitle=white,
  fonttitle=\bfseries,
  left=8pt, right=8pt, top=6pt, bottom=6pt,
  boxrule=0pt, leftrule=4pt,
  arc=4pt,
  title={#1}
}

% highlightbox — Yellow-bordered for notable findings
\newtcolorbox{highlightbox}[1][]{%
  colback=highlightyellow!8,
  colframe=highlightyellow,
  coltitle=white,
  fonttitle=\bfseries,
  left=8pt, right=8pt, top=6pt, bottom=6pt,
  boxrule=0pt, leftrule=4pt,
  arc=4pt,
  title={#1}
}

% assumptionbox — Gold full-border for assumptions
\newtcolorbox{assumptionbox}[1][]{%
  colback=primarygold!8,
  colframe=primarygold,
  coltitle=white,
  fonttitle=\bfseries,
  left=8pt, right=8pt, top=6pt, bottom=6pt,
  boxrule=2pt,
  arc=4pt,
  title={#1}
}

% resultbox — Gold-bordered with shadow for main results
\newtcolorbox{resultbox}[1][]{%
  colback=primarygold!15,
  colframe=primarygold,
  coltitle=white,
  fonttitle=\bfseries,
  left=8pt, right=8pt, top=6pt, bottom=6pt,
  boxrule=2pt,
  arc=4pt,
  shadow={1pt}{-1pt}{0pt}{black!12},
  title={#1}
}

% eqbox — Subtle blue background for equations
\newtcolorbox{eqbox}{%
  colback=primaryblue!4,
  colframe=primaryblue!4,
  left=8pt, right=8pt, top=4pt, bottom=4pt,
  boxrule=0pt,
  arc=0pt
}

% softbox — Subtle gold italic for remarks
\newtcolorbox{softbox}{%
  colback=primarygold!6,
  colframe=primarygold!6,
  left=8pt, right=8pt, top=4pt, bottom=4pt,
  boxrule=0pt,
  arc=0pt,
  fontupper=\itshape
}

% quotebox — Italic with decorative quote mark
\newtcolorbox{quotebox}{%
  colback=primarygold!8,
  colframe=primarygold!80!black,
  left=16pt, right=8pt, top=6pt, bottom=6pt,
  boxrule=0pt, leftrule=4pt,
  arc=0pt, outer arc=4pt,
  fontupper=\itshape,
  before upper={\textcolor{primarygold!40}{\Large``}\,}
}

% --- Text color commands ---
\newcommand{\hi}[1]{{\bfseries\color{primaryblue}#1}}
\newcommand{\higold}[1]{{\bfseries\color{primarygold}#1}}
\newcommand{\hiyellow}[1]{{\bfseries\color{highlightyellow}#1}}
\newcommand{\higreen}[1]{{\bfseries\color{positivegreen}#1}}
\newcommand{\hired}[1]{{\bfseries\color{negativered}#1}}
\newcommand{\positive}[1]{{\bfseries\color{positivegreen}#1}}
\newcommand{\negative}[1]{{\bfseries\color{negativered}#1}}
\newcommand{\neutral}[1]{{\color{gray}#1}}

% --- Bibliography ---
\usepackage[
  backend=bibtex,
  style=authoryear,
  natbib=true,
  maxcitenames=2
]{biblatex}

% --- Title page template (compact, no internal \vfill) ---
\setbeamertemplate{title page}{
  \begin{center}
    \vspace{1em}
    {\Large\bfseries\color{primaryblue}\inserttitle}

    \vspace{0.4em}
    {\normalsize\color{primarygold}\insertsubtitle}

    \vspace{0.8em}
    {\normalsize\insertauthor}

    \vspace{0.2em}
    {\small\textcolor{primarygold}{\insertinstitute}}
  \end{center}
}

% --- Frame title formatting ---
\setbeamerfont{frametitle}{size=\large, series=\bfseries}
\setbeamertemplate{frametitle}{%
  \vspace{0.3em}
  \insertframetitle
  \vspace{0.1em}
  {\color{primarygold!30}\hrule height 2pt}
}


\addbibresource{../Bibliography_base.bib}

\title{DreamDojo: Foundation World Model\\for Dexterous Robotics}
\subtitle{Gao, Ye, Yang, Zheng, Jang et al.\ (NVIDIA, 2026)}
\author{Yunsung Lee}
\institute{WoRV / MaumAI}
\date{}

\begin{document}

% ==============================================================================
% PART 1: MOTIVATION & BACKGROUND
% ==============================================================================

% --- Slide 1: Title ---
\begin{frame}[plain]
\maketitle
\vfill
\begin{center}
\includegraphics[width=0.65\textwidth,
  height=0.45\textheight,keepaspectratio]{%
  ../Figures/DreamDojo/teaser_compressed.pdf}
\end{center}
\end{frame}

% --- Slide 1b: Roadmap ---
\begin{frame}{Roadmap}
\begin{enumerate}
    \item \hi{Motivation} --- why robot data is not enough
    \item \hi{Latent Action Model} --- self-supervised actions
    \item \hi{Architecture \& Training} --- building the world model
    \item \hi{Distillation} --- making it real-time
    \item \hi{Results} --- benchmarks and applications
    \item \hi{Code Insights} --- under the hood
    \item \hi{Discussion} --- takeaways and landscape
\end{enumerate}
\end{frame}

% --- Slide 2: DreamZero Recap ---
\begin{frame}{DreamZero Recap: WAMs for Zero-shot Policies}
\hi{World Action Models (WAMs):}
\begin{itemize}
    \item \hi{Joint} video + action generation
          from a single DiT backbone
    \item Autoregressive chunk-wise inference
          with ground-truth injection
    \item \positive{2$\times$} over VLAs,
          \positive{38$\times$} speedup to real-time 7Hz
\end{itemize}

\vspace{0.5em}
\begin{keybox}[Key Paradigm]
WAMs generate ``\hi{how to move}'' by predicting
``\hi{what happens next}'' ---
inverse dynamics emerge from video understanding.
\end{keybox}
\end{frame}

% --- Slide 2b: From DreamZero to DreamDojo ---
\begin{frame}{From DreamZero to DreamDojo}
\begin{columns}[T]
\column{0.48\textwidth}
\hi{DreamZero (WAM = Policy):}
\begin{itemize}
    \item Joint video + action prediction
    \item \hi{Inverse} dynamics:
          video $\rightarrow$ action
    \item \positive{Zero-shot policies}
    \item Trained on robot data
\end{itemize}

\column{0.48\textwidth}
\hi{DreamDojo (World Model):}
\begin{itemize}
    \item Action-conditioned video prediction
    \item \hi{Forward} dynamics:
          action $\rightarrow$ video
    \item \positive{Open-world simulation}
    \item Pretrained on \higold{human video}
\end{itemize}
\end{columns}

\vspace{0.5em}
\begin{keybox}[The Shift]
Same team, complementary paradigms:
DreamZero is a \hi{policy}, DreamDojo is a \hi{simulator}.
\end{keybox}
\end{frame}

% --- Slide 3: The Data Gap ---
\begin{frame}{The Data Gap: Robot Data Is Not Enough}
\hi{The fundamental bottleneck:}
\begin{itemize}
    \item Existing robot datasets: limited coverage,
          narrow distributions
    \item DROID: 350 hours, 86 skills, 564 scenes
    \item AgiBot World: 2.9K hours, 87 skills, 106 scenes
    \item World models trained on robot data
          $\rightarrow$ \negative{in-distribution only}
\end{itemize}

\vspace{0.5em}
\begin{highlightbox}[The Core Problem]
Robot world models plateau at \negative{observed setups}
and are \negative{unresponsive to counterfactual actions}
--- actions the robot never performed in training
(e.g., patting a toy, reaching and missing).
Scaling teleoperation alone cannot cover the long tail.
\end{highlightbox}
\end{frame}

% --- Slide 4: The Human Video Insight ---
\begin{frame}{The Human Video Insight}
\begin{quotebox}
Despite the embodiment gap, the underlying physics
during interactions is largely consistent between
humans and robots, enabling effective knowledge transfer.
\end{quotebox}
\hfill --- \citet{gao2026dreamdojo}

\vspace{0.5em}
\begin{center}
\includegraphics[width=0.82\textwidth,
  height=0.40\textheight,keepaspectratio]{%
  ../Figures/DreamDojo/dataset_compressed.pdf}
\end{center}

\vspace{0.3em}
\hi{Key idea:} Human videos capture the same physics ---
contact, gravity, deformation, tool use ---
just with different end-effectors.
\end{frame}

% --- Slide 5: DreamDojo-HV at Scale ---
\begin{frame}{DreamDojo-HV: The Largest WM Dataset}
\begin{resultbox}[Dataset Scale]
\begin{itemize}
    \item \positive{44,711 hours} of egocentric video
    \item \positive{1,179K} trajectories
    \item \positive{$\sim$6,015} unique skills
    \item \positive{$>$1,135K} unique scenes
\end{itemize}
\end{resultbox}

\vspace{0.2em}
\hi{vs.\ prior largest (AgiBot World):}
\positive{15$\times$} duration,
\positive{96$\times$} skills,
\positive{2,000$\times$} scenes

\vspace{0.2em}
\begin{center}
\includegraphics[width=0.70\textwidth,
  height=0.22\textheight,keepaspectratio]{%
  ../Figures/DreamDojo/mecka_compressed.pdf}
\end{center}
\end{frame}

% --- Slide 6: Three Data Sources ---
\begin{frame}{Three Data Sources}
\begin{methodbox}[Data Mixture]
\begin{columns}[T]
\column{0.32\textwidth}
\hi{In-lab (55h)}\\
{\small Manus gloves,\\
35 skills, 1 scene}

\column{0.32\textwidth}
\hi{EgoDex (829h)}\\
{\small Apple Vision Pro,\\
194 skills, 5 scenes}

\column{0.32\textwidth}
\hi{HV (43.8Kh)}\\
{\small Crowdsourced,\\
6,015 skills, 1.1M scenes}
\end{columns}
\end{methodbox}

\vspace{0.3em}
\hi{Sampling ratio:} In-lab : EgoDex : HV $=$ \higold{1 : 2 : 10}

\vspace{0.3em}
\begin{columns}[c]
\column{0.48\textwidth}
\begin{center}
\includegraphics[width=0.82\textwidth,
  height=0.20\textheight,keepaspectratio]{%
  ../Figures/DreamDojo/egodex.png}\\
{\small EgoDex sample}
\end{center}
\column{0.48\textwidth}
\begin{center}
\includegraphics[width=0.82\textwidth,
  height=0.20\textheight,keepaspectratio]{%
  ../Figures/DreamDojo/inlab.png}\\
{\small In-lab PSNR curves}
\end{center}
\end{columns}
\end{frame}

% --- Slide 7: System Overview ---
\begin{frame}{System Overview: Three-Phase Pipeline}
\begin{center}
\includegraphics[width=0.90\textwidth,
  height=0.45\textheight,keepaspectratio]{%
  ../Figures/DreamDojo/overview_compressed.pdf}
\end{center}

\vspace{0.2em}
\hi{Phase 1:} Pretrain on human videos with \higold{latent actions}\\
\hi{Phase 2:} Post-train on target robot data (reset action layer)\\
\hi{Phase 3:} Distill to autoregressive real-time student

\end{frame}

% ==============================================================================
% PART 2: LATENT ACTION MODEL
% ==============================================================================

% --- Section Divider ---
\begin{frame}[plain]
\begin{center}
\vfill
{\Huge\hi{How Does DreamDojo Learn Actions?}}\\[1em]
{\large\higold{Latent Action Model}}
\vfill
\end{center}
\end{frame}

% --- Slide 8: The Action Label Problem ---
\begin{frame}{The Action Label Problem}
\hi{Human videos are vast but lack action labels.}

\vspace{0.5em}
\begin{itemize}
    \item Passive video prediction ignores \negative{causality}
          between observations and actions
    \item Off-the-shelf hand pose extractors (HaMeR)
          miss arm movements, locomotion,
          and fail under occlusion
    \item Converting heterogeneous action formats into
          a unified representation
          $\rightarrow$ significant engineering burden
\end{itemize}

\vspace{0.5em}
\begin{highlightbox}[The Dilemma]
Actionless pretraining wastes interaction knowledge.
Ground-truth actions require expensive capture devices.
We need a \hi{self-supervised} middle ground.
\end{highlightbox}
\end{frame}

% --- Slide 9: LAM Architecture ---
\begin{frame}{Latent Action Model (LAM)}
\begin{center}
\includegraphics[width=0.88\textwidth,
  height=0.38\textheight,keepaspectratio]{%
  ../Figures/DreamDojo/lam_compressed.pdf}
\end{center}

\vspace{0.2em}
\begin{eqbox}
\hi{VAE with information bottleneck:}
\begin{equation*}
\mathcal{L}^{pred}_{\theta,\phi}(f^{t+1})
  = \mathbb{E}_{q_{\phi}(\hat{a}|f^{t:t+1})}
    \log p_{\theta}(f^{t+1}|\hat{a},f^{t})
  - \beta\,D_{KL}(q_{\phi}(\hat{a}|f^{t:t+1})||p(\hat{a}))
\end{equation*}
{\small $\hat{a} \in \mathbb{R}^{32}$, \ \ $\beta=10^{-6}$
(tight bottleneck $\rightarrow$ forces action disentanglement)}
\end{eqbox}
\end{frame}

% --- Slide 10: Latent Actions Cross Embodiments ---
\begin{frame}{Latent Actions Cross Embodiments}
\begin{keybox}[Cross-Embodiment Transfer]
The same latent action $\hat{a}$ captures the
\hi{same physical interaction} regardless of whether
it is performed by a human hand, a gripper,
or a dexterous robot.
\end{keybox}

\vspace{0.5em}
\hi{LAM Architecture (700M params):}
\begin{itemize}
    \item Spatiotemporal Transformer
          \citep{bruce2024genie}
    \item 24 encoder blocks $+$ 24 decoder blocks
    \item Encoder: $(f^t, f^{t+1}) \rightarrow \hat{a}$
          \ \ Decoder: $(\hat{a}, f^t) \rightarrow \hat{f}^{t+1}$
    \item Trained on mixture of \hi{all} datasets
          (human + robot)
\end{itemize}

\end{frame}

% --- Slide 11: Latent Actions vs Alternatives ---
\begin{frame}{Latent Actions vs.\ Alternatives}
\begin{center}
{\small
\begin{tabular}{lccl}
\toprule
\textbf{Method} & \textbf{PSNR$\uparrow$}
  & \textbf{LPIPS$\downarrow$} & \textbf{Note}\\
\midrule
w/o pretrain & 20.576 & 0.222 & No human data \\
Action-free & 20.797 & 0.222 & Passive video \\
\hi{Latent action} & \positive{20.913}
  & \positive{0.219} & Self-supervised \\
\midrule
\neutral{Retargeted GT} & \neutral{20.960}
  & \neutral{0.219} & \neutral{Requires gloves} \\
\bottomrule
\end{tabular}
}
\end{center}

\vspace{0.3em}
\begin{methodbox}[Key Finding]
Latent actions achieve \positive{near-parity}
with ground-truth action labels ---
without any capture devices.
The gap: $<$0.05 PSNR, 0.000 LPIPS.
\end{methodbox}

\vspace{0.3em}
\begin{softbox}
\higold{My Take:} This is the enabler for 44K-hour scale.
You cannot put Manus gloves on a million crowdworkers,
but you \hi{can} run a VAE on their videos.
\end{softbox}
\end{frame}

% ==============================================================================
% PART 3: ARCHITECTURE & TRAINING
% ==============================================================================

% --- Section Divider ---
\begin{frame}[plain]
\begin{center}
\vfill
{\Huge\hi{Building the World Model}}\\[1em]
{\large\higold{Architecture \& Training}}
\vfill
\end{center}
\end{frame}

% --- Slide 12: Cosmos-Predict2.5 Backbone ---
\begin{frame}{Cosmos-Predict2.5 Backbone}
\begin{methodbox}[Foundation]
\begin{itemize}
    \item Latent video diffusion model
          \citep{ali2025world}
    \item \hi{WAN 2.2 tokenizer}:
          $4\times$ temporal compression\\
          {\small (video latent $x^i$ corresponds to
          4 pixel frames $f^{i:i+4}$)}
    \item \hi{DiT} blocks \citep{peebles2023scalable}
          with cross-attention for text,
          AdaLN for timestep
    \item Trained with \hi{flow matching} loss
    \item Two variants: \higold{2B} and \higold{14B} parameters
\end{itemize}
\end{methodbox}

\vspace{0.3em}
\hi{Training setup:}
\begin{itemize}
    \item 140K steps pretrain (256 H100s),
          50K steps post-train (128 H100s)
    \item Resolution: $640 \times 480$,
          sequences of 13 latent frames
    \item EMA maintained throughout
\end{itemize}
\end{frame}

% --- Slide 12b: Flow Matching in 60 Seconds ---
\begin{frame}{Flow Matching in 60 Seconds}
\begin{keybox}[Intuition]
Learn a \hi{velocity field} that transports
noise $\rightarrow$ data along straight paths.
Simpler and faster than standard diffusion.
\end{keybox}

\vspace{0.5em}
\begin{eqbox}
\hi{Flow matching loss:}
\begin{equation*}
\mathcal{L}_{\text{flow}}(\theta)
  = \mathbb{E}_{\mathbf{x},\epsilon,\mathbf{c},t}
    \left\|\mathbf{u}(\mathbf{x}_{t},t,\mathbf{c};\theta)
    - \mathbf{v}_{t}\right\|^2
\end{equation*}
{\small where $\mathbf{v}_t = \epsilon - \mathbf{x}$
(velocity = noise $-$ clean data)}
\end{eqbox}

\vspace{0.5em}
\begin{itemize}
    \item $t=0$: pure noise.
          $t=1$: clean data.
    \item Model predicts the \hi{velocity} pointing
          from noise toward data
    \item Conditions $\mathbf{c}$: text, frames,
          and \higold{actions} (for world models)
\end{itemize}
\end{frame}

% --- Slide 13: Two Key Innovations ---
\begin{frame}{Two Key Innovations for Controllability}
\begin{center}
\includegraphics[width=0.80\textwidth,
  height=0.35\textheight,keepaspectratio]{%
  ../Figures/DreamDojo/design_compressed.pdf}
\end{center}

\begin{methodbox}[Design Choices]
\begin{enumerate}
    \item \hi{Relative actions:} Rebaseline every 4 timesteps\\
          {\small Concentrates action space
          $\rightarrow$ better generalization to compositions}
    \item \hi{Chunked injection:} 4 consecutive actions
          $a^{t:t+4}$ sent to matching latent frame\\
          {\small Respects causality
          $\rightarrow$ future actions do not leak}
\end{enumerate}
\end{methodbox}
\end{frame}

% --- Slide 14: Action Embedding Equation ---
\begin{frame}{Action Embedding: AdaLN Injection}
\begin{eqbox}
\hi{Action injection pathway:}
\begin{equation*}
\text{actions} \xrightarrow{\text{MLP}}
  \text{action embedding}
  \xrightarrow{+}
  \text{timestep embedding}
  \xrightarrow{\text{AdaLN}}
  \text{DiT blocks}
\end{equation*}
{\small Scale, shift, and gate modulations
in every DiT block}
\end{eqbox}

\vspace{0.5em}
\hi{Key implementation details:}
\begin{itemize}
    \item Last layer of action MLP
          \hi{zero-initialized}
          \citep{zhang2023adding}\\
          {\small Avoids perturbing pretrained state
          at start of training}
    \item During post-training: \hi{reset first layer only},
          finetune all weights
    \item Classifier-free guidance \hi{disabled}
          (empirically no benefit)
\end{itemize}
\end{frame}

% --- Slide 14b: Action Embedding Code Pattern ---
\begin{frame}{Action Embedding: Code Insights}
\begin{methodbox}[Implementation Pattern]
\begin{itemize}
    \item Latent actions projected via lightweight MLP
          to match timestep embedding dim
    \item \hi{Zero-init} on last MLP layer
          $\rightarrow$ identity-like at init
    \item Added to timestep embeddings before AdaLN
    \item Dual pathway: actions modulate through
          \hi{both} timestep AND scale/shift/gate
\end{itemize}
\end{methodbox}

\vspace{0.5em}
\begin{softbox}
\higold{My Take:} The dual-path action injection is more
sophisticated than the paper describes --- actions
modulate DiT through \hi{both timestep AND AdaLN}
scale/shift/gate pathways. The zero-init trick from
ControlNet is crucial for preserving pretrained physics.
\end{softbox}
\end{frame}

% --- Slide 15: Temporal Consistency Loss ---
\begin{frame}{Temporal Consistency Loss}
\hi{Problem:} Flow matching supervises frames individually,
ignoring temporal correlations.

\vspace{0.3em}
\begin{eqbox}
\hi{Temporal consistency loss:}
\begin{equation*}
\mathcal{L}_{\text{temporal}}(\theta)
  = \mathbb{E}\Big[\sum^{K-1}_{i=1}
    \left\|(z^{i+1}-z^{i})-(v^{i+1}-v^{i})\right\|^2\Big]
\end{equation*}
\end{eqbox}

\vspace{0.2em}
\begin{eqbox}
\hi{Final objective:}
$\mathcal{L}_{\text{final}}
  = \mathcal{L}_{\text{flow}}
  + \lambda\,\mathcal{L}_{\text{temporal}},
  \ \ \lambda = 0.1$
\end{eqbox}

\end{frame}

% --- Slide 16: Post-Training ---
\begin{frame}{Post-Training: Adapting to Target Robots}
\begin{center}
\includegraphics[width=0.80\textwidth,
  height=0.34\textheight,keepaspectratio]{%
  ../Figures/DreamDojo/pretrain_compressed.pdf}
\end{center}

\vspace{0.2em}
\begin{methodbox}[Post-Training Recipe]
\begin{enumerate}
    \item \hi{Reset first layer} of action MLP
          (new action space)
    \item \hi{Finetune all weights},
          target data at $\sim$10 Hz
    \item 128 H100s, 50K steps, batch 512
\end{enumerate}
\end{methodbox}

\vspace{0.2em}
{\small Embodiments: GR-1, G1, AgiBot, YAM
--- each post-trained separately.
CFG disabled (no benefit).}
\end{frame}

% ==============================================================================
% PART 4: DISTILLATION
% ==============================================================================

% --- Section Divider ---
\begin{frame}[plain]
\begin{center}
\vfill
{\Huge\hi{Making It Real-Time}}\\[1em]
{\large\higold{Distillation Pipeline}}
\vfill
\end{center}
\end{frame}

% --- Slide 18: The Distillation Challenge ---
\begin{frame}{The Distillation Challenge}
\hi{Why distill?}
\begin{itemize}
    \item Teacher: \negative{35 denoising steps},
          bidirectional attention, 2.72 FPS
    \item Too slow for live teleoperation
          and online planning
\end{itemize}

\vspace{0.5em}
\begin{highlightbox}[Two Conversions Required]
\begin{enumerate}
    \item \hi{Bidirectional $\rightarrow$ Autoregressive}
          (causal attention, streaming)
    \item \hi{35 steps $\rightarrow$ 4 steps}
          (few-step generation)
\end{enumerate}
\end{highlightbox}

\vspace{0.3em}
Student initialized from teacher weights,
with bidirectional attention replaced by
\hi{causal attention} over a sliding window of 12 frames.
\end{frame}

% --- Slide 19a: Self Forcing Warmup ---
\begin{frame}{Self Forcing: Warmup Stage}
\hi{Stage 1: Warmup} --- student regresses to teacher ODE solutions

\vspace{0.5em}
\begin{eqbox}
\begin{equation*}
\mathcal{L}_{\text{warmup}}
  = \mathbb{E}_{x,t}\|G_{\text{student}}(x_{t},t)-x_{0}\|^2
\end{equation*}
{\small Student generates via teacher forcing (clean context)}
\end{eqbox}

\vspace{0.5em}
\hi{How it works:}
\begin{itemize}
    \item Student receives \hi{clean} context frames
          (teacher forcing)
    \item Learns to map noisy inputs to clean data
          in fewer steps
    \item Provides a stable initialization
          before self-generated training
\end{itemize}
\end{frame}

% --- Slide 19b: Self Forcing Distillation ---
\begin{frame}{Self Forcing: Distillation Stage}
\hi{Stage 2: Distillation} --- student trains on its own outputs
\citep{huang2025self}

\vspace{0.3em}
\begin{eqbox}
\begin{equation*}
\nabla\mathcal{L}_{\text{distill}}
  = -\mathbb{E}_{z,t}\Big[
    (s_{\text{real}}(x_{t},t)
    - s_{\text{fake}}(x_{t},t))
    \frac{dG_{\text{student}}}{d\theta}\Big]
\end{equation*}
{\small KL divergence between teacher and student
distributions \citep{yin2024one}}
\end{eqbox}

\vspace{0.2em}
{\small $s_{\text{real}}$: score from real (teacher-generated) data
\quad $s_{\text{fake}}$: score from fake (student-generated) data}

\vspace{0.3em}
\hi{Extended rollout:} Student generates $N'>N$ frames,
loss on last $N$ --- reduces long-horizon drift.
\end{frame}

% --- Slide 20: Distillation Results ---
\begin{frame}{Distillation Results}
\begin{resultbox}[Speed vs.\ Quality]
\begin{center}
{\small
\begin{tabular}{lcccccc}
\toprule
& \textbf{PSNR$\uparrow$} & \textbf{SSIM$\uparrow$}
  & \textbf{LPIPS$\downarrow$}
  & \textbf{FPS$\uparrow$} & \textbf{Pred}
  & \textbf{Ctx}\\
\midrule
Teacher & 14.09 & 0.442 & 0.412
  & 2.72 & 12 & 1\\
Student & 13.15 & 0.379 & 0.485
  & \positive{$\sim$10} & 4 & 12\\
\bottomrule
\end{tabular}
}
\end{center}
\end{resultbox}

\vspace{0.3em}
\begin{columns}[T]
\column{0.48\textwidth}
\begin{center}
\includegraphics[width=\textwidth,
  height=0.28\textheight,keepaspectratio]{%
  ../Figures/DreamDojo/long_compressed.pdf}\\
{\small 1-minute rollouts}
\end{center}
\column{0.48\textwidth}
\begin{center}
\includegraphics[width=\textwidth,
  height=0.28\textheight,keepaspectratio]{%
  ../Figures/DreamDojo/context_compressed.pdf}\\
{\small Context advantage}
\end{center}
\end{columns}
\end{frame}

% ==============================================================================
% PART 5: RESULTS
% ==============================================================================

% --- Section Divider ---
\begin{frame}[plain]
\begin{center}
\vfill
{\Huge\hi{Does It Actually Work?}}\\[1em]
{\large\higold{Experimental Results}}
\vfill
\end{center}
\end{frame}

% --- Slide 21: Benchmarks ---
\begin{frame}{Benchmark Construction: 6 OOD Eval Sets}
\begin{center}
\includegraphics[width=0.85\textwidth,
  height=0.35\textheight,keepaspectratio]{%
  ../Figures/DreamDojo/benchmark_compressed.pdf}
\end{center}

\begin{methodbox}[6 Evaluation Sets on GR-1]
{\small
\hi{1.} In-lab Eval \
\hi{2.} EgoDex Eval \
\hi{3.} HV Eval \
\hi{4.} Counterfactual Eval\\
\hi{5.} EgoDex-novel (Gemini-edited BG) \
\hi{6.} HV-novel (Gemini-edited BG)\\[0.3em]
Metrics: PSNR, SSIM, LPIPS (auto);
human preference for novel scenes.
}
\end{methodbox}
\end{frame}

% --- Slide 22: Data Scaling Works ---
\begin{frame}{Data Scaling: More Diversity $=$ Better OOD}
\begin{center}
{\small
\begin{tabular}{lccc}
\toprule
\textbf{Data Mixture} & \textbf{In-lab}
  & \textbf{EgoDex} & \textbf{HV}\\
& PSNR$\uparrow$ & PSNR$\uparrow$ & PSNR$\uparrow$\\
\midrule
No pretrain (Cosmos) & 20.58 & 19.95 & 18.27\\
In-lab only & 20.91 & 20.27 & 18.62\\
+ EgoDex & 20.97 & 20.33 & 18.71\\
+ DreamDojo-HV & 21.02 & \positive{20.41} & 18.72\\
\midrule
DreamDojo-2B & 21.11 & 20.41 & 18.81\\
DreamDojo-14B & \positive{21.41} & \positive{20.53}
  & \positive{18.92}\\
\bottomrule
\end{tabular}
}
\end{center}

\vspace{0.3em}
\begin{keybox}[Monotonic Improvement]
Adding more human data \positive{consistently improves}
all OOD benchmarks --- both physics modeling
and counterfactual action following.
\end{keybox}
\end{frame}

% --- Slide 23: Human Preference ---
\begin{frame}{Human Preference: 73.5\% Win Rate}
\begin{resultbox}[Human Evaluation (12 volunteers)]
\begin{center}
{\small
\begin{tabular}{lcc}
\toprule
\textbf{Comparison} & \textbf{Physics}
  & \textbf{Action}\\
\midrule
DreamDojo-2B $>$ Cosmos
  & 62.5\% & 63.5\%\\
DreamDojo-14B $>$ Cosmos
  & \positive{73.5\%} & \positive{72.6\%}\\
DreamDojo-14B $>$ DreamDojo-2B
  & 72.5\% & 65.5\%\\
\bottomrule
\end{tabular}
}
\end{center}
\end{resultbox}

\vspace{0.5em}
\hi{Two axes evaluated:}
\begin{itemize}
    \item \hi{Physics correctness:} object permanence,
          shape consistency, contact causality
    \item \hi{Action following:} robot pose accuracy
          vs.\ ground-truth
\end{itemize}
\end{frame}

% --- Slide 24: Ablation ---
\begin{frame}{Ablation: What Matters Most?}
\begin{center}
{\small
\begin{tabular}{ccc|ccc}
\toprule
\textbf{Rel.} & \textbf{Chunk} & \textbf{Temp.}
  & \multicolumn{3}{c}{\textbf{Counterfactual Eval}}\\
& & & PSNR$\uparrow$ & SSIM$\uparrow$ & LPIPS$\downarrow$\\
\midrule
 & & & 19.45 & 0.768 & 0.211\\
\checkmark & & & 19.48 & 0.772 & 0.212\\
\checkmark & \checkmark &
  & \positive{20.78} & \positive{0.790}
  & \positive{0.193}\\
\checkmark & \checkmark & \checkmark
  & \positive{20.98} & \positive{0.796}
  & \positive{0.189}\\
\bottomrule
\end{tabular}
}
\end{center}

\vspace{0.3em}
\begin{keybox}[Takeaway]
\hi{Chunked injection} is the largest single gain
($+$1.30 PSNR). Relative actions alone have minimal effect,
but temporal loss adds a further $+$0.20 improvement.
\end{keybox}
\end{frame}

% --- Slide 25: Multi-Embodiment ---
\begin{frame}{Multi-Embodiment Support}
\begin{keybox}[Four Robot Platforms]
Post-trained on GR-1, G1, AgiBot, and YAM
--- each with separate post-training runs
using the same pretrained foundation.
\end{keybox}

\vspace{0.5em}
\hi{Per-embodiment adaptation:}
\begin{itemize}
    \item Reset action MLP first layer for new action space
    \item Finetune all weights (50K steps)
    \item Same pretrained physics transfers across all platforms
\end{itemize}

\end{frame}

% --- Slide 26: Policy Evaluation ---
\begin{frame}{Policy Evaluation: A Reliable Simulator}
\begin{columns}[T]
\column{0.38\textwidth}
\begin{center}
\includegraphics[width=\textwidth,
  height=0.45\textheight,keepaspectratio]{%
  ../Figures/DreamDojo/correlation_compressed.pdf}
\end{center}

\column{0.58\textwidth}
\begin{resultbox}[Correlation with Reality]
\begin{itemize}
    \item Pearson $r$ = \positive{0.995}
    \item MMRV = \positive{0.003}
\end{itemize}
\end{resultbox}

\vspace{0.3em}
\hi{Protocol:}
\begin{itemize}
    \item AgiBot fruit packing (20 scenes)
    \item GR00T N1.5 policy checkpoints
    \item $\sim$80s rollouts in real world
          \emph{and} DreamDojo
    \item Success = fruits packed / 5
\end{itemize}
\end{columns}

\vspace{0.3em}
{\small DreamDojo can reliably rank policies without
deploying on a real robot.}
\end{frame}

% --- Slide 27: MPC & Teleoperation ---
\begin{frame}{MPC \& Live Teleoperation}
\begin{columns}[T]
\column{0.48\textwidth}
\hi{Model-Based Planning:}
\begin{center}
\includegraphics[width=\textwidth,
  height=0.30\textheight,keepaspectratio]{%
  ../Figures/DreamDojo/mpc_compressed.pdf}
\end{center}
\begin{itemize}
    \item 5 policy checkpoints $\rightarrow$ action proposals
    \item DreamDojo predicts futures, value model selects
    \item \positive{+17\%} success rate over best checkpoint
\end{itemize}

\column{0.48\textwidth}
\hi{Live Teleoperation:}
\begin{center}
\includegraphics[width=\textwidth,
  height=0.30\textheight,keepaspectratio]{%
  ../Figures/DreamDojo/teleop_compressed.pdf}
\end{center}
\begin{itemize}
    \item PICO VR controller $\rightarrow$ G1 actions
    \item Distilled model on RTX 5090
    \item Real-time interactive simulation
\end{itemize}
\end{columns}

\vspace{0.2em}
\begin{keybox}
World models as \hi{simulators} unlock planning,
evaluation, and teleoperation ---
all from a single foundation.
\end{keybox}
\end{frame}

% ==============================================================================
% PART 6: CODE INSIGHTS
% ==============================================================================

% --- Section Divider ---
\begin{frame}[plain]
\begin{center}
\vfill
{\Huge\hi{Under the Hood}}\\[1em]
{\large\higold{Code-Level Insights}}
\vfill
\end{center}
\end{frame}

% --- Slide 27b: Why Look at Code? ---
\begin{frame}{Why Look at the Code?}
\begin{keybox}[Production-Grade NVIDIA Code]
DreamDojo's codebase reveals engineering decisions
the paper glosses over --- these details matter
for reproducibility and understanding.
\end{keybox}

\vspace{0.5em}
\hi{What we can learn:}
\begin{itemize}
    \item How the LAM bottleneck actually works
    \item How actions modulate the DiT
    \item How three different data loaders
          handle heterogeneous sources
    \item How relative actions use proper SE(3) math
\end{itemize}
\end{frame}

% --- Slide 28: Code LAM ---
\begin{frame}{Code: Latent Action Model}
\begin{methodbox}[Implementation Details (Not in Paper)]
\begin{itemize}
    \item Class \texttt{LatentActionModel} wraps
          \hi{asymmetric} encoder--decoder
    \item Encoder: \texttt{SpatioTemporalTransformer}
          (joint space-time attention across frames)
    \item Decoder: \texttt{SpatioTransformer}
          (spatial-only --- no temporal mixing)
    \item Action prompt initialized as:\\
          {\small \texttt{nn.Parameter(torch.empty(1, 1, 1, patch\_token\_dim))}}
\end{itemize}
\end{methodbox}

\vspace{0.3em}
\hi{Key asymmetry:} The encoder sees \hi{both} frames
jointly (spatiotemporal), while the decoder reconstructs
from action + \hi{single} frame (spatial only).
This forces all temporal information through the
32-dim bottleneck.
\end{frame}

% --- Slide 29: Code DiT ---
\begin{frame}{Code: DiT Action Injection}
\begin{methodbox}[Concrete Code Patterns]
\begin{itemize}
    \item Temporal rearrange:\\
          {\small \texttt{rearrange(action, "b 1 (t d) -> b t d", t=...)}}
    \item \hi{Two separate MLPs} for action embedding:\\
          {\small \texttt{action\_embedder\_B\_D}: global timestep path\\
          \texttt{action\_embedder\_B\_3D}: per-frame spatial path}
    \item Each MLP last layer \hi{zero-initialized}
    \item Embeddings added to timestep \emph{before} AdaLN
\end{itemize}
\end{methodbox}

\vspace{0.3em}
\begin{softbox}
\higold{My Take:} Two separate action MLPs ---
not one --- feed into different pathways.
This dual-path design is more sophisticated
than the paper describes and explains why
chunked injection has such a large ablation effect.
\end{softbox}
\end{frame}

% --- Slide 30: Code Pipeline ---
\begin{frame}{Code: Data Pipeline}
\begin{methodbox}[Three Distinct Loaders]
\begin{itemize}
    \item \hi{MANO loader:} For EgoDex/In-lab
          with hand pose annotations
    \item \hi{VideoOnly loader:} For DreamDojo-HV
          (latent actions extracted online)
    \item \hi{LeRobot loader:} For target robot data
          with ground-truth joint actions
\end{itemize}
\end{methodbox}

\vspace{0.3em}
\hi{Dual-resolution processing:}\\
{\small LAM operates at $320\times240$ (efficiency),
while the world model trains at $640\times480$ (quality).
Separate tokenizer passes for each resolution.}
\end{frame}

% --- Slide 31: Code Relative Actions ---
\begin{frame}{Code: Relative Action Computation}
\begin{methodbox}[SE(3) Math for Relative Actions]
\begin{itemize}
    \item Translation: body-frame deltas,
          not naive world-frame subtraction
    \item Rotation: proper rotation matrix computation
          for relative orientation
    \item Scale factor: $\times$20 applied
          to normalize action magnitudes
    \item Rebaseline every 4 timesteps
          (matching tokenizer compression)
\end{itemize}
\end{methodbox}

\vspace{0.5em}
\begin{softbox}
\higold{My Take:} The relative action computation uses
proper SE(3) math --- translation deltas in body frame,
not naive subtraction. The \hi{20$\times$ scaling factor}
is a practical detail the paper omits entirely.
These ``boring'' engineering choices are
what make the model actually work.
\end{softbox}
\end{frame}

% ==============================================================================
% PART 7: DISCUSSION
% ==============================================================================

% --- Section Divider ---
\begin{frame}[plain]
\begin{center}
\vfill
{\Huge\hi{What Does It All Mean?}}\\[1em]
{\large\higold{Discussion \& Takeaways}}
\vfill
\end{center}
\end{frame}

% --- Slide 32: Limitations ---
\begin{frame}{Limitations \& Honest Assessment}
\begin{highlightbox}[Paper-Acknowledged Limitations]
\begin{itemize}
    \item \negative{Uncommon actions} (slapping, fast waving)
          still fail
    \item Policy eval \negative{overestimates}
          success rates vs.\ real world
    \item No multi-view simulation support
    \item Post-training knowledge retention not studied
\end{itemize}
\end{highlightbox}

\vspace{0.3em}
\begin{softbox}
\higold{My Take:} The proprietary 43.8K-hour dataset
is both the biggest strength and weakness ---
it makes the results hard to reproduce.
The eval overestimation suggests the model
struggles with nuanced failure modes.
\end{softbox}
\end{frame}

% --- Slide 33: DreamDojo vs DreamZero Paradigms ---
\begin{frame}{DreamDojo vs.\ DreamZero: Two Paradigms}
\begin{columns}[T]
\column{0.48\textwidth}
\hi{DreamZero (WAM = Policy):}
\begin{itemize}
    \item \hi{Inverse} dynamics
    \item Video + action \emph{jointly} generated
    \item Zero-shot policy transfer
    \item Trained on robot data only
    \item 500 hours robot teleoperation
\end{itemize}

\column{0.48\textwidth}
\hi{DreamDojo (World Model):}
\begin{itemize}
    \item \hi{Forward} dynamics
    \item Action $\rightarrow$ video prediction
    \item Simulation, planning, teleoperation
    \item Pretrained on \higold{44K hours human video}
    \item Adapted to robots via post-training
\end{itemize}
\end{columns}

\vspace{0.3em}
\begin{keybox}[Complementary, Not Competing]
DreamZero asks ``what action produces this future?''\\
DreamDojo asks ``what future does this action produce?''
\end{keybox}
\end{frame}

% --- Slide 33b: DreamDojo vs DreamZero Technical ---
\begin{frame}{DreamDojo vs.\ DreamZero: Technical Comparison}
\begin{center}
{\small
\begin{tabular}{lcc}
\toprule
& \textbf{DreamZero} & \textbf{DreamDojo}\\
\midrule
Dynamics & Inverse (IDM) & Forward (WM)\\
Backbone & WAN 2.1 & Cosmos-Predict2.5\\
Tokenizer & WAN 2.1 ($4\times$) & WAN 2.2 ($4\times$)\\
Conditioning & Joint video+action & Action $\rightarrow$ video\\
Data & Robot only & Human + robot\\
Scale & 500h robot & 44.7Kh total\\
Latent actions & No & Yes (32-dim VAE)\\
Distillation & Flash (1-step) & Self Forcing (4-step)\\
Speed & 7Hz (2$\times$GB200) & $\sim$10 FPS (1$\times$H100)\\
\bottomrule
\end{tabular}
}
\end{center}

\vspace{0.3em}
\begin{softbox}
\higold{My Take:} The Cosmos-Predict2.5 backbone gives
DreamDojo stronger video priors, while the latent action
framework enables human video pretraining at scale.
Different tools for different jobs.
\end{softbox}
\end{frame}

% --- Slide 34: Landscape ---
\begin{frame}{The World Model Landscape}
\begin{softbox}
\hi{DreamDojo} is part of a rapid convergence in
world models for robotics:
\end{softbox}

\vspace{0.3em}
\begin{itemize}
    \item \hi{AdaWorld} \citep{gao2025adaworld}:
          latent actions for adaptable WMs (tabletop)
    \item \hi{Genie} \citep{bruce2024genie}:
          interactive environments from video
          (discrete control)
    \item \hi{DexWM} \citep{goswami2025dexwm}:
          human videos for dexterous manipulation
          (smaller scale)
    \item \hi{DreamDojo}: first foundation WM at
          \positive{44K hours} with
          \positive{continuous} latent actions
\end{itemize}

\vspace{0.3em}
\hi{The trend:} Scale up human video,
bridge the embodiment gap with latent actions,
distill for real-time interaction.
\end{frame}

% --- Slide 35: Key Takeaways ---
\begin{frame}{Key Takeaways}
\begin{enumerate}
    \item \hi{Human video pretraining works}
          --- 44K hours of ego video
          $\rightarrow$ physics that generalizes
    \item \hi{Latent actions bridge embodiments}
          --- self-supervised, near GT quality,
          infinitely scalable
    \item \hi{Chunked + relative actions}
          are the biggest controllability gains
    \item \hi{Self Forcing distillation}
          enables real-time interaction ($\sim$10 FPS)
    \item \hi{Policy evaluation} correlates
          with reality ($r=0.995$)
    \item \higold{World models are becoming
          general-purpose simulators}
          for planning, eval, and teleoperation
\end{enumerate}

\vspace{0.3em}
\begin{center}
\hi{The future of robot world models is
human-video-native.}
\end{center}
\end{frame}

% --- Slide 36: Thank You ---
\begin{frame}{Thank You \& References}
\begin{columns}[T]
\column{0.48\textwidth}
\hi{Paper \& Code:}
\begin{itemize}
    \item \citet{gao2026dreamdojo}
    \item arXiv: 2602.06949
    \item Code: \texttt{github.com/NVIDIA/DreamDojo}
\end{itemize}

\vspace{0.3em}
\hi{DreamZero (predecessor):}
\begin{itemize}
    \item \citet{ye2026dreamzero}
    \item \texttt{dreamzero0.github.io}
\end{itemize}

\column{0.48\textwidth}
\hi{Key References:}
\begin{itemize}
    \item Cosmos-Predict2.5 \citep{ali2025world}
    \item Self Forcing \citep{huang2025self}
    \item AdaWorld \citep{gao2025adaworld}
    \item Genie \citep{bruce2024genie}
    \item DiT \citep{peebles2023scalable}
    \item Flow matching \citep{lipman2022flow}
\end{itemize}
\end{columns}

\vspace{0.5em}
\begin{center}
\Large\hi{Q\&A}
\end{center}
\end{frame}

% --- Bibliography ---
\begin{frame}[allowframebreaks]{References}
\printbibliography[heading=none]
\end{frame}

\end{document}
